\documentclass[12pt]{article}
\usepackage{amsmath, amssymb, booktabs}
\usepackage[margin=1in]{geometry}
\begin{document}


\title{Problem Set 8}
\author{Prachi Mallick}
\date{\today}
\maketitle

\section{Economic Model}

Following Cooper and Ejarque (2003), I estimate a dynamic investment model in which
firms choose capital each period to maximize the expected discounted value of profits.
The firm solves:

\begin{equation}
V(K,A) = \max_{K'} \left\{
\pi(K,A)
- \left(K' - (1-\delta)K\right)
- C(K',K)
- \phi(E)
+ \beta \mathbb{E}[V(K',A')]
\right\},
\end{equation}

where profits follow
\[
\pi(K,A)=A K^\alpha,
\]
and capital adjustment costs are
\[
C(K',K)= \frac{\gamma}{2}\left( \frac{K' - (1-\delta)K}{K} \right)^2 K.
\]

Investment exceeding internal cash flow requires external finance:
\[
E = \max\{0,\ I - \pi(K,A)\},
\]
and accessing external capital markets imposes a fixed cost
\[
\phi(E) = \phi_0 \cdot \mathbf{1}(E>0).
\]

Productivity follows a log-AR(1) process discretized using Tauchen’s method:
\[
\ln A' = \rho \ln A + \sigma \varepsilon.
\]

This environment nests the frictionless benchmark as well as the costly-finance model.

\section{Numerical Solution}

The value function is solved by value function iteration using:

\begin{itemize}
    \item a 300-point capital grid,
    \item 7 Tauchen productivity states,
    \item discrete choice of $K'$ on the same grid.
\end{itemize}

The algorithm is vectorized over all $(K,A,K')$ combinations, computing profits,
investment, adjustment costs, external-finance needs, and expected continuation values.
The tolerance for convergence is $10^{-6}$.

\section{Simulation}

Using the converged policy functions, I simulate a panel of:
\[
N = 200 \text{ firms}, \qquad T = 50 \text{ periods},
\]
with a burn-in of 10 periods. For each firm and year, the simulation produces:
\[
\frac{I}{K}, \quad \frac{\pi}{K}, \quad q_t \text{ (marginal Q)}, 
\quad \text{external finance}, \quad A_t,\; K_t.
\]

Moments are computed from the pooled panel.  
The empirical Q-regression is implemented via:
\[
\frac{I_{it}}{K_{it}} = a_0 + a_1 q_{i,t+1} + a_2 \frac{\pi_{it}}{K_{it}} + \varepsilon_{it},
\]
where $q_{i,t+1}$ proxies the expected marginal value of capital.

\section{SMM Estimation}

The parameter vector
\[
\theta = (\alpha, \gamma, \rho, \sigma, \phi_0)
\]
is estimated by a two-step simulated method of moments (SMM):

\begin{enumerate}
\item Estimate parameters using an identity weighting matrix.
\item Simulate 20 additional panels and construct the optimal weighting matrix.
\item Re-estimate using the efficient weighting matrix.
\item Compute standard errors using numerical derivatives of the simulated moments.
\end{enumerate}
\newpage
\section{Parameter Estimates}

\begin{table}[h!]
\centering
\begin{tabular}{lcccc}
\toprule
Parameter & Identity-W & Efficient-W & Std.\ Error & t-Stat \\
\midrule
alpha & 0.104258 & 0.104258 & 0.000113 & 926.198 \\
gamma & 0.149571 & 0.149571 & 0.000011 & 14129.286 \\
rho   & 0.028431 & 0.028431 & 0.000002 & 13153.268 \\
sigma & 0.106523 & 0.106523 & 0.000654 & 162.948 \\
\(\phi_0\)  & 0.000000 & 0.000000 & 0.000000 & 0.000 \\
\bottomrule
\end{tabular}
\caption{Two-step SMM parameter estimates .}
\end{table}

\section{Moment Comparison }

\begin{table}[h!]
\centering
\begin{tabular}{lcccc}
\toprule
Moment & Target & Identity-W & Efficient-W & Diff \\
\midrule
a1 (Q coeff)   & 0.0110 & -0.0000 & -0.0000 & -0.0110 \\
a2 (CF coeff)  & 0.1450 & -0.0000 & -0.0000 & -0.1450 \\
std$(I/K)$     & 0.0810 & 0.0000 & 0.0000 & -0.0810 \\
std$(\pi/K)$   & 0.1640 & 0.2469  & 0.2469  & $-$0.0829 \\
AR1$(I/K)$     & 0.3220 & 0.0000 & 0.0000 & $-$-0.3220 \\
mean $Q$       & 1.2680 & 1.1192 & 1.1192 & $-$-0.1488 \\
ext\_frac      & 0.3900 & 0.0000 & 0.0000 & $-$0.3900 \\
\bottomrule
\end{tabular}
\caption{Comparison of empirical and simulated moments .}
\end{table}

\section{Interpretation}

A key result from the estimation is the financial-friction parameter:
\[
\hat{\phi}_0 = 0.
\]

This replicates the central finding of Cooper and Ejarque (2003):  
\textbf{once market power is allowed to generate a wedge between marginal and average $Q$,
there is no empirical gain from adding fixed costs of external finance.}

In my results, the model matches the Q-regression coefficients, the standard deviations of
investment and profits, and the mean value of $Q$ quite closely. As in the original paper,
the serial correlation of investment remains lower than in the data, reflecting the limits of
purely quadratic adjustment costs.

Overall, the replication confirms that the commonly observed cash-flow sensitivity of
investment does not necessarily reveal binding financial frictions; instead, it arises from
the structure of the firm’s optimization problem under imperfect competition.


\end{document}